\documentclass[a4paper]{article}

\usepackage[utf8]{inputenc}
\usepackage[a4paper, margin=3.5cm]{geometry}

\usepackage{graphicx}
\usepackage{url}
\usepackage[spanish,es-tabla]{babel}
\usepackage{fancyhdr}
\usepackage{float}
\usepackage[breaklinks=true]{hyperref}
\usepackage[newfloat]{minted}
\usepackage[nolist]{acronym}
\usepackage{glossaries}
\usepackage[
    type={CC},
    modifier={by-nc-sa},
    version={3.0},
]{doclicense}

\hypersetup{breaklinks=true}

\renewcommand{\headrulewidth}{0.6pt}
\renewcommand{\footrulewidth}{0.6pt}

\pagestyle{fancy}
\setlength\headheight{50pt}
\lhead{\includegraphics[height=1.5cm]{Practicas/logos/upm_logo}}
\chead{Métodos generativos\\\vspace{.5em} Práctica 1\\\vspace{-.1em}}
\rhead{\includegraphics[height=1.5cm]{Practicas/logos/etsisi_logo}}
\lfoot{\textbf{Tema 2:} Arquitecturas de métodos generativos}
\cfoot{}
\rfoot{\thepage}

\parskip 1.1ex % paragraph spacing
\newglossaryentry{AE}{name={AE},description={Autoencoder}}
\newglossaryentry{VAE}{name={VAE},description={Variational Autoencoder}}
\newglossaryentry{GAN}{name={GAN},description={Generative Adversarial Network}}


\begin{document}

\section*{Normativa}

\begin{itemize}
    \item La práctica se realizará en grupos de 5 alumnos con la siguiente distribución:
    \begin{itemize}
        \item 2 alumnos del Grado de Ingeniería del Software
        \item 1 alumno del Grado de Ingeniería de Computadores
        \item 1 alumno del Grado de Ciencia de Datos e Inteligencia Artificial
        \item 1 alumno comodín de cualquier titulación
    \end{itemize}
    \item Si se detecta cualquier sospecha de copia conllevará un 0.
    \item No se podrán utilizar modelos entrenados con la técnica de \textit{fine-tunning} ni arquitecturas distintas a las vistas en clase.
    \item Se penalizará el uso de LLMs sin criterio para redactar el contenido o generar el código. El alumno siempre es el último responsable de su trabajo.
\end{itemize}

\section{Objetivo de la práctica}
\textbf{Dataset:} Para la realización de la práctica y la presentación de los resultados se va a utilizar un dataset homogéneo de caras de humanos, para poder crear nuevas ''personas'' a nuestro antojo. \href{https://github.com/switchablenorms/CelebAMask-HQ}{\textbf{CelebAMask-HQ}} es un dataset de 30.000 imágenes a color de 512x512 píxeles que contiene imágenes de personas, así como máscaras de segmentación para 19 de sus atributos faciales.

\begin{figure}[H]
    \centering
    \includegraphics[width=0.5\linewidth]{Practicas/2627/Practica_1/CelebAMask-HQ.jpeg}
\end{figure}

Partiendo de este dataset se ha generado un subconjunto más pequeño, de 64x64 píxeles para la realización de esta práctica, que está disponible en Kaggle (previa invitación) \href{https://www.kaggle.com/competitions/mg-practica-1-generando-caras}{\nolinkurl{https://www.kaggle.com/competitions/mg-practica-1-generando-caras}}. Durante esta práctica sólo serán necesarias las imágenes de las caras, no las máscaras de segmentación.

\section{Evaluación}
La evaluación de la práctica consistirá en una competición de Kaggle junto a una breve presentación en clase. Todos los grupos deberán participar en ambas o de lo contrario la práctica se considerará no entregada.

\textbf{Competición de Kaggle:} Utilizando cada uno de los modelos de redes neuronales generativos que se lista a continuación y el dataset citado anteriormente, se debe entrenar un modelo capaz de generar imágenes parecidas al dataset para cada arquitectura.
\begin{itemize}
    \item \gls{AE}: entrenamiento e inferencia
    \item \gls{VAE}: entrenamiento e inferencia
\end{itemize}

Elige uno de los modelos utilizados en el apartado anterior para participar en la competición entre todos los alumnos de este curso. Para ello, vamos a utilizar las herramientas que nos ofrece la plataforma de Kaggle (URL de invitación a la competición: \href{https://www.kaggle.com/t/313d806078e46a21725d1e3ce786c35f}{\textbf{pincha aqui}}), en el enlace encontrareis la competición entre vosotros y las puntuaciones obtenidas en cada subida. El punto de este apartado lo conseguirán solo los 3 mejores resultados al finalizar la competición.

\textbf{Presentación:} Cada grupo deberá realizar una breve presentación de 6 minutos utilizando un póster o diapositivas complementarias, donde se explique el trabajo realizado y se haga un análisis crítico de los resultados obtenidos en la competición de Kaggle.

La presentación debe incluir:
\begin{itemize}
    \item La justificación del diseño de las arquitecturas utilizadas y sus entrenamientos. También se deben incluir reflexiones sobre pruebas realizadas para obtener los mejores resultados.
    \item Un análisis de los resultados obtenidos visualmente.
    \item Explicación de los resultados subidos a Moodle.
    \item Generación de imágenes que fusionen características de distintas personas que \textbf{no pertenezcan al dataset}. Por ejemplo, podéis usar fotos de los profesores de la asignatura\footnote{Os lo ponemos fácil, aqui os pasamos alguna: \href{https://master-sostec.etsisi.upm.es/wp-content/uploads/2019/03/edgar.jpg}{\nolinkurl{https://master-sostec.etsisi.upm.es/wp-content/uploads/2019/03/edgar.jpg}} y \href{https://blogs.upm.es/gietema/wp-content/uploads/sites/1082/2025/10/Guille-886x1024.jpeg}{\nolinkurl{https://blogs.upm.es/gietema/wp-content/uploads/sites/1082/2025/10/Guille-886x1024.jpeg}}} para generar una cara que fusione los rasgos faciales de ambos.
\end{itemize}

\section{Entrega en Moodle}
Además se deberán subir los notebooks correspondientes para los entrenos realizados, como prueba del funcionamiento de lo anteriormente descrito.

Asimismo, para poder ser evaluada, la entrega deberá llevar asociada una participación en la competición de Kaggle previamente descrita.

\doclicenseThis
\end{document}